\documentclass[10pt,a4paper]{article}

\usepackage[utf8]{inputenc}
\usepackage[T1]{fontenc}
\usepackage{lmodern}
\usepackage[a4paper,top=1.8cm,bottom=1.8cm,left=1.8cm,right=1.8cm]{geometry}
\usepackage{xcolor}
\usepackage{tikz}
\usetikzlibrary{calc,shapes.geometric,backgrounds}
\usepackage{booktabs}
\usepackage{tabularx}
\usepackage{enumitem}
\usepackage{titlesec}
\usepackage{fancyhdr}
\usepackage{microtype}
\usepackage[colorlinks=true,linkcolor=primary,urlcolor=secondary]{hyperref}

% =========================================================
% Color Palette
% =========================================================
\definecolor{primary}{RGB}{30, 58, 138}        % Deep navy blue
\definecolor{secondary}{RGB}{2, 132, 199}      % Cyan/Sky blue
\definecolor{darkslate}{RGB}{51, 65, 85}       % Slate grey
\definecolor{bordergray}{RGB}{209, 213, 219}   % Light grey border
\definecolor{boxbg}{RGB}{248, 250, 252}        % Light subtle card bg

% Terminal Colors (Light Grey Theme)
\definecolor{termbg}{RGB}{246, 248, 250}       % Light grey background
\definecolor{termbar}{RGB}{229, 231, 235}      % Slightly darker grey title bar
\definecolor{termborder}{RGB}{209, 213, 219}   % Border grey
\definecolor{termtext}{RGB}{31, 41, 55}        % Dark charcoal text
\definecolor{termuser}{RGB}{22, 101, 52}       % Dark green for user
\definecolor{termhost}{RGB}{21, 128, 61}       % Green for host
\definecolor{termpath}{RGB}{29, 78, 216}       % Blue for directory path
\definecolor{termcmd}{RGB}{17, 24, 39}         % Bold black for commands
\definecolor{termout}{RGB}{55, 65, 81}         % Clear dark grey for outputs
\definecolor{termdir}{RGB}{37, 99, 235}        % Blue for directories
\definecolor{termfile}{RGB}{31, 41, 55}        % Dark for regular files
\definecolor{termhidden}{RGB}{107, 114, 128}   % Muted grey for hidden files
\definecolor{termsize}{RGB}{180, 83, 9}        % Amber for sizes
\definecolor{termperm}{RGB}{79, 70, 229}       % Indigo for permissions
\definecolor{termkey}{RGB}{194, 65, 12}        % Rust orange for editor shortcut keys

% =========================================================
% Section Headings (Clean, No Underline)
% =========================================================
\titleformat{\section}
  {\color{primary}\large\bfseries\sffamily}
  {\thesection.}{0.5em}{}
\titlespacing*{\section}{0pt}{8pt}{3pt}

\titleformat{\subsection}
  {\color{secondary}\normalsize\bfseries\sffamily}
  {\thesubsection}{0.5em}{}
\titlespacing*{\subsection}{0pt}{6pt}{2pt}

% =========================================================
% Headers and Footers (Completely Clean Footers)
% =========================================================
\pagestyle{fancy}
\fancyhf{}
\renewcommand{\headrulewidth}{0.4pt}
\renewcommand{\footrulewidth}{0pt}
\fancyhead[L]{\small\sffamily\color{darkslate} Operating Systems Lab}
\fancyhead[R]{\small\sffamily\color{darkslate} Lab Report: Viewing \& Editing Files}
\fancyfoot{}

% First page style: no header at top, no footer at bottom
\fancypagestyle{firstpage}{%
  \fancyhf{}
  \renewcommand{\headrulewidth}{0pt}
  \renewcommand{\footrulewidth}{0pt}
}

% =========================================================
% Terminal Window Macro (Light Grey with Crisp Edges)
% =========================================================
\newcommand{\termprompt}[1][\textasciitilde]{%
  \textbf{\textcolor{termuser}{ratul}}\textcolor{termtext}{@}\textbf{\textcolor{termhost}{firefly}}\textcolor{termtext}{:}\textbf{\textcolor{termpath}{#1}}\textcolor{termcmd}{\$}\ %
}

\newcommand{\terminalwindow}[3][\linewidth]{%
\begin{center}
\begin{tikzpicture}
  % 1. Terminal text on main layer
  \node[
      text=termtext,
      font=\ttfamily\fontsize{7.1}{9.0}\selectfont,
      align=left,
      text width=\dimexpr#1-0.95cm\relax,
      inner xsep=7pt,
      inner ysep=7pt,
      anchor=north west
  ] (body) at (0.32cm, -0.56cm) {#3};

  % Coordinates
  \coordinate (topleft) at (0, 0);
  \coordinate (bottomright) at ($(body.south -| #1, 0) + (0, -6pt)$);
  \coordinate (titlecenter) at ($0.5*(topleft) + 0.5*(#1, 0) + (0, -0.28cm)$);

  % 2. Background and Title bar on background layer
  \begin{scope}[on background layer]
    % Light grey terminal body
    \fill[termbg, rounded corners=6pt] (topleft) rectangle (bottomright);

    % Title bar (clipped to match rounded top corners)
    \begin{scope}
      \clip[rounded corners=6pt] (topleft) rectangle (bottomright);
      \fill[termbar] (topleft) rectangle (#1, -0.56cm);
      \draw[termborder, line width=0.5pt] (0, -0.56cm) -- (#1, -0.56cm);
    \end{scope}

    % Continuous outer border (smooth, no distortion)
    \draw[termborder, line width=0.6pt, rounded corners=6pt] (topleft) rectangle (bottomright);
  \end{scope}

  % 3. Window control dots
  \fill[red!80!black] (0.36cm, -0.28cm) circle (0.07cm);
  \fill[yellow!80!black] (0.63cm, -0.28cm) circle (0.07cm);
  \fill[green!75!black] (0.90cm, -0.28cm) circle (0.07cm);

  % 4. Centered terminal title
  \node[
      text=gray!40!black,
      font=\sffamily\bfseries\fontsize{7.2}{8}\selectfont,
      anchor=center
  ] at (titlecenter) {#2};
\end{tikzpicture}
\end{center}
\vspace{-6pt}
}

\newcommand{\cmd}[1]{\texttt{\textbf{#1}}}
\newcommand{\key}[1]{\textbf{\textcolor{termkey}{#1}}}

\begin{document}
\thispagestyle{firstpage}

% =========================================================
% Title: Only Experiment Title at the Top
% =========================================================
\begin{center}
  {\Large\bfseries\sffamily\color{primary} Experiment Title: Viewing and Editing Files in Linux Terminal}\\[3pt]
  {\normalsize\sffamily\textbf{Topic:} File Viewing Tools (\cmd{cat}, \cmd{less}, \cmd{head}, \cmd{tail}) and Text Editors (\cmd{nano}, \cmd{vim})}
\end{center}
\vspace{-4pt}

% =========================================================
% 1. Objective
% =========================================================
\section{Objective}
The objective of this lab is to learn how to view file contents quickly from the command line and edit configuration or source files using terminal-based text editors. By completing this experiment, we aim to:
\begin{itemize}[noitemsep,topsep=1pt,leftmargin=1.2em]
  \item Inspect whole text files, file headers, and file tails using \cmd{cat}, \cmd{less}, \cmd{head}, and \cmd{tail}.
  \item Check software presence using \cmd{which} and install missing editors using \cmd{sudo apt install}.
  \item Master basic text editing with \cmd{nano}: saving, exiting, cut, copy, paste, undo, redo, searching, and replacing.
  \item Master modal text editing with \cmd{vim}: switching modes, saving, quitting, yank (copy), put (paste), delete (cut), undo, redo, searching, and substitute (replace).
\end{itemize}

% =========================================================
% 2. Equipment / Software Requirements
% =========================================================
\section{Equipment / Software Requirements}
\vspace{-2pt}
\begin{table}[h!]
\centering\footnotesize
\renewcommand{\arraystretch}{1.08}
\begin{tabularx}{\linewidth}{lXlX}
\toprule
\textbf{Item} & \textbf{Details} & \textbf{Item} & \textbf{Details} \\
\midrule
\textbf{Operating System} & Linux (Ubuntu / Debian / NixOS) & \textbf{Host Machine} & \texttt{firefly} \\
\textbf{User Account} & \texttt{ratul} (with sudo privileges) & \textbf{Terminal Shell} & GNU Bash 5.2 \\
\textbf{Viewing Utilities} & GNU Coreutils (\texttt{cat}, \texttt{head}, \texttt{tail}), \texttt{less} & \textbf{Text Editors} & GNU Nano 7.2, Vim 9.1 \\
\bottomrule
\end{tabularx}
\end{table}
\vspace{-8pt}

% =========================================================
% 3. Theoretical Background
% =========================================================
\section{Theoretical Background}
In Linux, administrators and programmers frequently need to read logs and modify system files without a graphical user interface (GUI). Tools in Linux fall into two groups: \textbf{pagers/viewers} (for reading without modifying) and \textbf{terminal text editors} (for creating and editing).

\subsection{File Viewing Commands}
\vspace{-2pt}
\begin{table}[h!]
\centering\footnotesize
\renewcommand{\arraystretch}{1.06}
\begin{tabularx}{\linewidth}{llp{5.4cm}X}
\toprule
\textbf{Command} & \textbf{Syntax} & \textbf{Simple Meaning / Purpose} & \textbf{Example} \\
\midrule
\cmd{cat} & \texttt{cat <file>} & Dumps the entire file onto the terminal screen. & \texttt{cat config.txt} \\
\cmd{less} & \texttt{less <file>} & Opens file in a scrollable viewer (page by page). & \texttt{less system.log} \\
\cmd{head} & \texttt{head -n <N> <file>} & Shows the first $N$ lines of a file (default is 10 lines). & \texttt{head -n 5 sample.txt} \\
\cmd{tail} & \texttt{tail -n <N> <file>} & Shows the last $N$ lines of a file (useful for recent logs). & \texttt{tail -n 3 server.log} \\
\bottomrule
\end{tabularx}
\end{table}
\vspace{-6pt}

\subsection{Comparing Nano and Vim}
\vspace{-2pt}
\begin{table}[h!]
\centering\footnotesize
\renewcommand{\arraystretch}{1.08}
\begin{tabularx}{\linewidth}{p{2.6cm}Xp{5.8cm}}
\toprule
\textbf{Editor} & \textbf{Design Style} & \textbf{Best Used For} \\
\midrule
\textbf{GNU Nano} & Simple, modeless editor. Always ready to type. Shortcut hints are shown at the bottom of the screen with \key{\^{}} (Ctrl) and \key{M-} (Alt). & Quick edits, configuration changes, and beginners who want immediate typing like Notepad. \\
\textbf{Vim} & Modal editor with distinct operational modes (\textit{Normal}, \textit{Insert}, and \textit{Command-line}). Keys act as commands in Normal mode. & Advanced editing, code writing, fast navigation without touching the mouse, available on almost every Unix server. \\
\bottomrule
\end{tabularx}
\end{table}

\newpage

% =========================================================
% 4. Quick Shortcut Reference Tables
% =========================================================
\section{Editor Commands and Shortcut Reference}

\subsection{GNU Nano Shortcuts}
In Nano, \key{\^{}} stands for the \key{Ctrl} key, and \key{M-} (Meta) stands for the \key{Alt} key:
\vspace{-2pt}
\begin{table}[h!]
\centering\footnotesize
\renewcommand{\arraystretch}{1.06}
\begin{tabularx}{\linewidth}{lp{3.5cm}X}
\toprule
\textbf{Action} & \textbf{Nano Shortcut / Command} & \textbf{Simple Explanation} \\
\midrule
\textbf{Check Presence} & \texttt{which nano} & Shows the location of nano executable on disk. \\
\textbf{Install} & \texttt{sudo apt install nano} & Installs nano package on Debian/Ubuntu systems. \\
\textbf{Save (WriteOut)} & \key{Ctrl + O} & Prompts for file name and saves changes to disk. \\
\textbf{Exit} & \key{Ctrl + X} & Closes nano (asks to save if changes were made). \\
\textbf{Cut Text} & \key{Ctrl + K} & Cuts (kills) the current line or selected text into the buffer. \\
\textbf{Copy Text} & \key{Alt + 6} (or \key{M-6}) & Copies the selected text or current line into the buffer. \\
\textbf{Paste Text} & \key{Ctrl + U} & Pastes (uncuts) whatever was last cut or copied. \\
\textbf{Undo} & \key{Alt + U} & Reverts the previous editing action. \\
\textbf{Redo} & \key{Alt + E} & Re-applies the undone editing action. \\
\textbf{Search (Find)} & \key{Ctrl + W} & Opens "Where Is" search prompt; press Enter to jump to match. \\
\textbf{Search \& Replace} & \key{Ctrl + \textbackslash} & Prompts for search word, then replacement word, and confirms each. \\
\bottomrule
\end{tabularx}
\end{table}
\vspace{-4pt}

\subsection{Vim Modal Commands}
In Vim, press \key{Esc} to ensure you are in \textbf{Normal Mode} before running commands:
\vspace{-2pt}
\begin{table}[h!]
\centering\footnotesize
\renewcommand{\arraystretch}{1.06}
\begin{tabularx}{\linewidth}{lp{3.5cm}X}
\toprule
\textbf{Action} & \textbf{Vim Command / Key} & \textbf{Simple Explanation} \\
\midrule
\textbf{Check Presence} & \texttt{which vim} & Displays the binary path for vim. \\
\textbf{Install} & \texttt{sudo apt install vim} & Installs full vim editor on Debian/Ubuntu systems. \\
\textbf{Enter Edit Mode} & \key{i} & Enters Insert Mode so typed keys appear as text. \\
\textbf{Save} & \texttt{:w} & Writes (saves) changes without quitting. \\
\textbf{Exit} & \texttt{:q} (or \texttt{:wq} / \texttt{:q!}) & Quits vim (\texttt{:wq} saves and quits; \texttt{:q!} quits without saving). \\
\textbf{Cut Line} & \key{dd} & Deletes (cuts) the whole current line into the paste register. \\
\textbf{Copy Line} & \key{yy} & Yanks (copies) the current line into memory. \\
\textbf{Paste} & \key{p} & Puts (pastes) copied/cut text below the current cursor line. \\
\textbf{Undo} & \key{u} & Undoes the last change made. \\
\textbf{Redo} & \key{Ctrl + r} & Redoes the undone change. \\
\textbf{Search Forward} & \texttt{/word} & Searches for "word"; press \key{n} for next match, \key{N} for previous. \\
\textbf{Search \& Replace} & \texttt{:\%s/old/new/g} & Replaces all occurrences of "old" with "new" across the entire file. \\
\bottomrule
\end{tabularx}
\end{table}
\vspace{-6pt}

% =========================================================
% 5. Procedure / Steps
% =========================================================
\section{Procedure / Steps}
We practiced viewing and editing through the following clear, numbered steps:
\begin{enumerate}[noitemsep,topsep=1pt,leftmargin=1.4em]
  \item Created a multi-line sample log file (\cmd{server.log}) with timestamps and error messages.
  \item Displayed the entire file using \cmd{cat server.log}.
  \item Extracted the top 3 lines using \cmd{head -n 3} and bottom 3 lines using \cmd{tail -n 3}.
  \item Scrolled through the file interactively with \cmd{less server.log} and exited with \key{q}.
  \item Checked nano installation with \cmd{which nano} (and simulated package update with \cmd{sudo apt install nano}).
  \item Opened nano, typed text, searched with \key{Ctrl+W}, replaced with \key{Ctrl+\textbackslash}, saved with \key{Ctrl+O}, and exited with \key{Ctrl+X}.
  \item Checked vim with \cmd{which vim}, opened a file, switched between Normal and Insert modes, performed copy (\key{yy}), paste (\key{p}), cut (\key{dd}), undo (\key{u}), redo (\key{Ctrl+r}), search (\cmd{/text}), replace (\cmd{:\%s/old/new/g}), and saved with \cmd{:wq}.
\end{enumerate}

\newpage

% =========================================================
% 6. Result / Output Demonstration
% =========================================================
\section{Result / Output Demonstration}
All commands were executed under user \texttt{ratul} on host machine \texttt{firefly}.

\subsection{Part 1: Viewing Files with \cmd{cat}, \cmd{head}, \cmd{tail}, and \cmd{less}}

\terminalwindow{ratul@firefly: \textasciitilde/lab4\ \textemdash\ bash}{
% View full file
\termprompt[\textasciitilde/lab4]\textbf{\textcolor{termcmd}{cat server.log}}

\textcolor{termout}{[2026-10-08 09:00:01] INFO  System startup complete}

\textcolor{termout}{[2026-10-08 09:05:12] DEBUG Loading database connection pool}

\textcolor{termout}{[2026-10-08 09:12:44] WARN  High CPU usage detected on core 2}

\textcolor{termout}{[2026-10-08 09:20:10] ERROR Port 8080 timeout connecting to worker}

\textcolor{termout}{[2026-10-08 09:30:00] INFO  Scheduled backup process finished}

\vspace{0.12cm}
% View first 2 lines
\termprompt[\textasciitilde/lab4]\textbf{\textcolor{termcmd}{head -n 2 server.log}}

\textcolor{termout}{[2026-10-08 09:00:01] INFO  System startup complete}

\textcolor{termout}{[2026-10-08 09:05:12] DEBUG Loading database connection pool}

\vspace{0.12cm}
% View last 2 lines
\termprompt[\textasciitilde/lab4]\textbf{\textcolor{termcmd}{tail -n 2 server.log}}

\textcolor{termout}{[2026-10-08 09:20:10] ERROR Port 8080 timeout connecting to worker}

\textcolor{termout}{[2026-10-08 09:30:00] INFO  Scheduled backup process finished}

\vspace{0.12cm}
% Open interactive pager
\termprompt[\textasciitilde/lab4]\textbf{\textcolor{termcmd}{less server.log}}

\textcolor{termhidden}{server.log (lines 1-5/5) (Press q to quit, / to search)}
}

\paragraph{Explanation of Part 1 Results:}
\begin{itemize}[noitemsep,topsep=1pt,leftmargin=1.2em]
  \item \cmd{cat}: Printed all lines directly to the terminal. Great for short files, but long files scroll off screen.
  \item \cmd{head -n 2}: Filtered and printed only the top 2 lines of the log.
  \item \cmd{tail -n 2}: Filtered and printed only the bottom 2 lines, which shows the most recent system events.
  \item \cmd{less}: Opened the file in a clean interactive view where we can scroll and search without modifying the file.
\end{itemize}

\subsection{Part 2: Installing and Editing with GNU Nano}

\terminalwindow{ratul@firefly: \textasciitilde/lab4\ \textemdash\ bash}{
% Check which nano
\termprompt[\textasciitilde/lab4]\textbf{\textcolor{termcmd}{which nano}}

\textcolor{termout}{/usr/bin/nano}

\vspace{0.08cm}
% Simulate install
\termprompt[\textasciitilde/lab4]\textbf{\textcolor{termcmd}{sudo apt install nano}}

\textcolor{termout}{nano is already the newest version (7.2-1ubuntu1).}

\vspace{0.12cm}
% Nano visual session
\termprompt[\textasciitilde/lab4]\textbf{\textcolor{termcmd}{nano app\_config.txt}}

\vspace{0.06cm}
\textcolor{termdir}{\textbf{GNU nano 7.2\ \ \ \ \ \ \ \ \ \ \ \ \ \ \ \ \ \ \ \ \ \ File: app\_config.txt\ \ \ \ \ \ \ \ \ \ \ \ \ \ \ \ \ \ \ \ \ \ Modified}}

\textcolor{termfile}{server\_mode = production}

\textcolor{termfile}{max\_connections = 500}

\textcolor{termfile}{port = 9000}

\vspace{0.06cm}
\textcolor{termdir}{\textbf{[ Replace "development" with "production" completed: 1 replacement ]}}

\textcolor{termhidden}{\textbf{\^{}G} Get Help\ \ \ \textbf{\^{}O} WriteOut\ \ \ \textbf{\^{}W} Where Is\ \ \ \textbf{\^{}K} Cut\ \ \ \ \ \ \textbf{M-U} Undo\ \ \ \ \ \ \textbf{\^{}X} Exit}

\textcolor{termhidden}{\textbf{\^{}R} Read File\ \ \textbf{\^{}\textbackslash} Replace\ \ \ \ \textbf{\^{}\_} Go To Line\ \textbf{\^{}U} Paste\ \ \ \ \textbf{M-E} Redo\ \ \ \ \ \ \textbf{M-6} Copy}

\vspace{0.08cm}
\termprompt[\textasciitilde/lab4]\textbf{\textcolor{termcmd}{cat app\_config.txt}}

\textcolor{termout}{server\_mode = production\ \ |\ \ max\_connections = 500\ \ |\ \ port = 9000}
}

\paragraph{Explanation of Part 2 Results:}
\begin{itemize}[noitemsep,topsep=1pt,leftmargin=1.2em]
  \item \cmd{which nano} confirmed nano is located at \texttt{/usr/bin/nano}. \cmd{sudo apt install nano} verified the package is ready.
  \item Inside nano, we used \key{Ctrl+W} to search, \key{Ctrl+\textbackslash} to replace \texttt{development} with \texttt{production}, \key{Ctrl+K} / \key{Ctrl+U} to cut and paste lines, and \key{Alt+U} / \key{Alt+E} for undo and redo.
  \item We saved with \key{Ctrl+O} and exited cleanly back to bash using \key{Ctrl+X}.
\end{itemize}

\newpage

% =========================================================
% 6. Result / Output Demonstration (Part 3)
% =========================================================
\subsection{Part 3: Installing and Editing with Vim}

\terminalwindow{ratul@firefly: \textasciitilde/lab4\ \textemdash\ bash}{
% Check which vim
\termprompt[\textasciitilde/lab4]\textbf{\textcolor{termcmd}{which vim}}

\textcolor{termout}{/usr/bin/vim}

\vspace{0.08cm}
% Simulate install
\termprompt[\textasciitilde/lab4]\textbf{\textcolor{termcmd}{sudo apt install vim}}

\textcolor{termout}{vim is already the newest version (2:9.1.0-1ubuntu1).}

\vspace{0.12cm}
% Vim visual session
\termprompt[\textasciitilde/lab4]\textbf{\textcolor{termcmd}{vim script.sh}}

\vspace{0.06cm}
\textcolor{termdir}{\textbf{\textasciitilde\ \ \ \ \ \ \ \ \ \ \ \ \ \ \ \ \ \ \ \ \ \ \ \ \ \ \ \ \ \ VIM - Vi IMproved 9.1}}

\textcolor{termfile}{1\ \ \#!/usr/bin/env bash}

\textcolor{termfile}{2\ \ echo "Starting backup service..."}

\textcolor{termfile}{3\ \ echo "Status: service running successfully."}

\textcolor{termfile}{4\ \ echo "Completed backup."}

\vspace{0.06cm}
\textcolor{termhost}{\textbf{-- INSERT --}}\ \ \ \ \ \ \ \ \ \ \ \ \ \ \ \ \ \ \ \ \ \ \ \ \ \ \ \ \ \ \ \ \ \ \ \ \ \ \ \ \ \ \ \ \ \ \ \ \ \ \ \ \ \ \ \ \ \ \ \ \ \ \ \textcolor{termout}{4,1\ \ \ \ \ \ \ \ \ \ All}

\textcolor{termuser}{\textbf{:\%s/task/service/g}}\ \ \ \ \textcolor{termout}{2 substitutions on 2 lines}

\textcolor{termuser}{\textbf{:wq}}\ \ \ \ \ \ \ \ \ \ \ \ \ \ \ \ \ \ \ \ \ \ \ \ \ \ \textcolor{termout}{"script.sh" 4L, 114B written}

\vspace{0.08cm}
\termprompt[\textasciitilde/lab4]\textbf{\textcolor{termcmd}{cat script.sh}}

\textcolor{termout}{\#!/usr/bin/env bash}

\textcolor{termout}{echo "Starting backup service..."}

\textcolor{termout}{echo "Status: service running successfully."}

\textcolor{termout}{echo "Completed backup."}

\vspace{0.08cm}
\termprompt[\textasciitilde/lab4]\textbf{\textcolor{termcmd}{\rule{0pt}{0.9em}\textunderscore}}
}

\paragraph{Explanation of Part 3 Results:}
\begin{itemize}[noitemsep,topsep=1pt,leftmargin=1.2em]
  \item \textbf{Checking \& Installing:} \cmd{which vim} verified the path \texttt{/usr/bin/vim}. \cmd{sudo apt install vim} confirmed the package is installed.
  \item \textbf{Modes:} Started in \textit{Normal Mode}. Pressed \key{i} to enter \textit{Insert Mode} and add code lines. Pressed \key{Esc} to return to Normal Mode.
  \item \textbf{Line Operations:} Pressed \key{yy} to yank (copy) line 2 and \key{p} to paste it. Pressed \key{dd} to cut (delete) an unwanted line.
  \item \textbf{Undo and Redo:} Pressed \key{u} to undo the deletion, and \key{Ctrl+r} to redo it.
  \item \textbf{Search and Replace:} Searched with \cmd{/backup}. Replaced words using the substitution command \cmd{:\%s/task/service/g}.
  \item \textbf{Save and Exit:} Saved and quit safely using \cmd{:wq}.
\end{itemize}

% =========================================================
% 7. Conclusion
% =========================================================
\section{Conclusion}
In this lab, we developed confidence in viewing and editing files completely from the Linux terminal:
\begin{enumerate}[noitemsep,topsep=1pt,leftmargin=1.4em]
  \item \textbf{Quick Inspection:} We used \cmd{cat} for quick output, \cmd{head} and \cmd{tail} to focus on file beginnings and ends, and \cmd{less} to comfortably scroll through larger logs.
  \item \textbf{Beginner Friendly Editing with Nano:} We learned that Nano offers a straightforward, modeless editor with clear on-screen shortcut keys for saving (\key{Ctrl+O}), cutting (\key{Ctrl+K}), pasting (\key{Ctrl+U}), searching (\key{Ctrl+W}), and replacing (\key{Ctrl+\textbackslash}).
  \item \textbf{Powerful Modal Editing with Vim:} We learned Vim's modal design, switching between Normal and Insert modes, executing rapid line operations (\key{yy}, \key{dd}, \key{p}), undoing mistakes (\key{u}), redoing (\key{Ctrl+r}), replacing text with \cmd{:\%s}, and saving/quitting with \cmd{:wq}.
\end{enumerate}
These command-line skills are fundamental for troubleshooting systems and managing Linux environments effectively.

\end{document}
